\chapter{Polynomial Arithmetic}

\section{Polynomial Arithmetic}

\subsection{Adding and Subtracting Polynomials}
\textbf{Monomial} is an expression that could be a numeral, a variable, or the product of numerals and variables within certain restrictions.\\
\\
1. Must have whole number exponents.\\
2. Variables do not appear under simplified radical expressions.\\
3. Denominators do not contain variables.\\

\textbf{Examples:} $-10$, $a$, $5t^2$, $3/8x^2y^3$ \\
\\
\textbf{Binomial} is an expression that is the sum of two monomials.\\
\textbf{Trinomial} is an expression that is the sum of three monomials.\\
\textbf{Polynomial} is an expression that is a monomial or the sum of two or more monomials.\\
\\

\subsection{Multiplying Polynomials}
1. $x^ax^b = x^{a+b}$ \\Multiplying monomials with the same base you keep the base and add the powers.\\\\
2. $(x^a)^b = x^{ab}$\\ Keep the base and multiply the powers when finding powers of a base.\\\\
3. $(xy)^a = x^ay^a$ \\Raise each factor in the proudct to the power when finding the powers of a product\\\\

\subsection{Special Products of Binomials}
\textbf{Conjugates} are two binomials with the same two terms but opposite signs separating the terms. They are conjugates of each other.\\
\\Example: $3x + 2$ \& $3x - 2$\\

\textbf{Difference of Squares} is the \textit{product} of conjugates which produces a special pattern.\\
\\Example: $(x + y)(x - y) = x^2 - y^2$\\
\\
\textbf{Square Trinomial} is the pattern produced by squaring a binomial.\\
\\Example: $(x + y)^2 = x^2 + 2xy + y^2$ \& $(x - y)^2 = x^2 - 2xy + y^2$\\
\\
Notice how only the \textit{first} sign is changed when subtracting a squared binomial.
\\ \\

\subsection{Dividing Polynomials}
\begin{center}
	\framebox{$a^m/a^n = a^{m-n}$} \\
\end{center}
This works as long as a does not equal zero.\\
\\
\begin{center}
	\framebox{$a^-n = 1/a^n$ \& $1/a^-n = a^n$} \\
\end{center}
This works as long as \textit{a} does not equal zero.\\ Also works in reverse.
\\

\subsection{synthetic Division}
\textbf{Synthetic Division} is a shortcut for polynomial division when the divisor is the form of x-a. Only the coefficients are used when dividing.\\
\\
Example: Divide $(3x^2 + 2x - 11)$ by $(x-3)$\\
\\
First step would be multiplying $(x-3)$ by $3x^2$ giving you $3x^3 - 9x^2$. You subtract this and keep repeating until you are left with a remainder. The remainder is added at the end by being divided by the divisor.

